% Part: first-order-logic
% Chapter: tableaux
% Section: provability-quantifiers

\documentclass[../../../include/open-logic-section]{subfiles}

\begin{document}

\olfileid{fol}{tab}{qpr}
\olsection{\usetoken{S}{derivability} lan Kuantor}

\begin{explain}
  Teorema kelengkapan uga mbutuhake supaya aturan tableaux ngasilake
  kasunyatan ngenani~$\Proves$ kang ditetepake ing bagean iki.
\end{explain}

\begin{thm}
\ollabel{thm:strong-generalization} Yen $c$ iku konstanta kang ora
muncul ing $\Gamma$ utawa~$!A(x)$ lan $\Gamma \Proves !A(c)$, mula
$\Gamma \Proves \lforall[x][!A(x)]$.
\end{thm}

\begin{proof}
Upamana $\Gamma \Proves !A(c)$, yaiku ana $!B_1$, \dots, $!B_n
\in \Gamma$ lan !!{tableau} katutup kanggo
\begin{align*}
\{ \sFmla{\False}{!A(c)}, &
\sFmla{\True}{!B_1}, \dots, \sFmla{\True}{!B_n} \}.
\intertext{Kita kudu nuduhake yen uga ana !!{tableau} katutup kanggo}
\{ \sFmla{\False}{\lforall[x][!A(x)]}, &
\sFmla{\True}{!B_1}, \dots, \sFmla{\True}{!B_n} \}.
\end{align*}
Jupuken !!{tableau} katutup mau lan gantinen asumsi kapisan nganggo
$\sFmla{\False}{\lforall[x][!A(x)]}$, banjur lebokna
$\sFmla{\False}{!A(c)}$ sawise asumsi-asumsi.
\begin{center}
\begin{tableau}{not line numbering}
  [\sFmla{\False}{\formula{A}(c)}
    [\sFmla{\True}{\formula{B}_1}
      [\vdots
      [\sFmla{\True}{\formula{B}_n}, 
        [][][]]]]]
\end{tableau}\qquad
\begin{tableau}{not line numbering}
  [\sFmla{\False}{\lforall[x][\formula{A}(x)]}
    [\sFmla{\True}{\formula{B}_1}
      [\vdots
        [\sFmla{\True}{\formula{B}_n},
          [\sFmla{\False}{\formula{A}(c)}
        [][][]]]]]]
\end{tableau}
\end{center}
Tableau kasebut tetep katutup, amarga kabeh !!{sentence}s kang
sadurunge cumawis minangka asumsi isih cumawis ing pucuk
!!{tableau}. Larik kang dilebokake minangka asil aplikasi bener saka
$\TRule{\False}{\lforall}$, amarga !!{constant}~$c$ ora muncul ing
$!B_1$, \dots, $!B_n$ utawa $\lforall[x][!A(x)]$, yaiku ora muncul
ing ndhuwur larik kang dilebokake ing !!{tableau} anyar.
\end{proof}

\begin{prop}
\ollabel{prop:provability-quantifiers}
\begin{tagenumerate}{prvEx,prvAll}
\tagitem{prvEx}{$!A(t) \Proves \lexists[x][!A(x)]$.}{}
\tagitem{prvAll}{$\lforall[x][!A(x)] \Proves !A(t)$.}{}
\end{tagenumerate}
\end{prop}

\begin{proof}
\begin{tagenumerate}{prvEx,prvAll}
\tagitem{prvEx}{Iki !!{tableau} katutup kanggo
  $\sFmla{\False}{\lexists[x][!A(x)]}, \sFmla{\True}{!A(t)}$:
  \begin{oltableau}
    [\sFmla{\False}{\lexists[x][\formula{A}(x)]}, just = \TAss
      [\sFmla{\True}{\formula{A}(t)}, just = \TAss
        [\sFmla{\False}{\formula{A}(t)}, just = {\TRule{\False}{\lexists}[1]}, close]]]
    \end{oltableau}}{}
\tagitem{prvAll}{Iki !!{tableau} katutup kanggo
  $\sFmla{\False}{\formula{A}(t)}, \sFmla{\True}{\lforall[x][!A(x)]}, $:
  \begin{oltableau}
    [\sFmla{\False}{\formula{A}(t)}, just = \TAss
      [\sFmla{\True}{\lforall[x][\formula{A}(x)]}, just = \TAss
        [\sFmla{\True}{\formula{A}(t)}, just = {\TRule{\True}{\lforall}[2]}, close]]]
    \end{oltableau}}{}
\end{tagenumerate}
\end{proof}

\end{document}
